% \iffalse meta-comment
%
% hit-date.dtx 是日期格式
%
% \fi
%
% \section{日期格式}
%
% \changes{v3.2a}{2026/08/29}{日期字段收 ISO 数字，各处按语言排、写到哪级排到哪级，非 ISO
%   原样排并报警告；\cs{today} 一族取东八区的今天}
%
% 封面与声明页上的日期有几种排法：\texttt{2026 年 8 月}、\texttt{2026 年 8 月 7 日}、
% \texttt{二〇二六年八月}、\texttt{二〇二六年八月七日}。这里把四种都定义好，
% 配置里按变体挑一种。
%
% \cs{CJKtoday}\oarg{0\textbar 1\textbar 2} 换 \cs{today} 的排法：\texttt{0} 是
% 基础类原本的，\texttt{1} 是阿拉伯数字，\texttt{2} 是汉字数字。默认 \texttt{1}。
%
%    \begin{macrocode}
%<*shared-date>
\cs_set:Npn \CJK@todaysmall@short
  {
    \int_use:N \g__hit_engine_today_year_int 年~
    \int_use:N \g__hit_engine_today_month_int 月
  }
\cs_set:Npn \CJK@todaysmall
  {
    \int_use:N \g__hit_engine_today_year_int 年~
    \int_use:N \g__hit_engine_today_month_int 月~
    \int_use:N \g__hit_engine_today_day_int 日
  }
\cs_set:Npn \CJK@todaybig@short
  {
    \zhdigits { \int_use:N \g__hit_engine_today_year_int } 年
    \zhnumber { \int_use:N \g__hit_engine_today_month_int } 月
  }
\cs_set:Npn \CJK@todaybig
  {
    \zhdigits { \int_use:N \g__hit_engine_today_year_int } 年
    \zhnumber { \int_use:N \g__hit_engine_today_month_int } 月
    \zhnumber { \int_use:N \g__hit_engine_today_day_int } 日
  }
\cs_set:Npn \CJK@today { \CJK@todaysmall }
\NewDocumentCommand \CJKtoday { O{1} }
  {
    \int_case:nnF {#1}
      {
        { 0 } { \cs_set:Npn \CJK@today { \CJK@todaysave } }
        { 1 } { \cs_set:Npn \CJK@today { \CJK@todaysmall } }
        { 2 } { \cs_set:Npn \CJK@today { \CJK@todaybig } }
      }
      { }
  }
\cs_new:Npn \hit@month@en
  {
    \int_case:nn { \g__hit_engine_today_month_int }
      {
        { 1 } { January }   { 2 } { February } { 3 }  { March }
        { 4 } { April }     { 5 } { May }      { 6 }  { June }
        { 7 } { July }      { 8 } { August }   { 9 }  { September }
        { 10 } { October }  { 11 } { November } { 12 } { December }
      }
  }
%    \end{macrocode}
%
% \subsection{ISO 日期}
%
% \textbf{日期字段填数字，排法交给模板。} 封面上中文写「2026 年 6 月」、英文写
% 「June, 2026」，同一个日期两处写法不同。让填表的人各写一遍容易写岔，所以字段
% 一律收 \texttt{yyyy}、\texttt{yyyy-mm} 或 \texttt{yyyy-mm-dd}，用到的地方按
% 自己的语言排出来。
%
% \textbf{写到哪级排到哪级。} 从前每个字段声明一个精度：写细了静默截断（写
% \texttt{2026-06-05} 印出「2026 年 6 月」，一声不响丢一级），写粗了报一句催
% 补齐。查过四处落点，规范一处也没要求到日——本科范例封面印「2022 年 6 月」，
% 旁边批注的原话是「\emph{年、月}用阿拉伯数字」；本科内封「答辩日期：2022 年
% 6 月」；研究生封面与内封都是年月。从前本科声明成到日，与范例和批注都对不上。
% 所以精度交还给写的人：写年月排年月，写年月日排年月日，不截不催。各处的默认
% 精度落在\emph{默认取值}上——模板给的今天是年月（见
% \file{src/config/hit-terms.dtx}），深圳本科的开题中期封面按 \issue{220}
% 单独给到日。
%
% 两种情形：
%
% \begin{longtblr}[caption={日期字段的写法与结果}, label={tab:impl-date}]
%   {colspec={l l}}
% \toprule
% 写的 & 结果 \\ \midrule
% 合 ISO & 按用处的语言、写到的那一级排 \\
% 不合 ISO & 原样排，另报一句让填表的人自己负责 \\
% \bottomrule
% \end{longtblr}
%    \begin{macrocode}
\tl_new:N \l__hit_date_year_tl
\tl_new:N \l__hit_date_month_tl
\tl_new:N \l__hit_date_day_tl
\bool_new:N \l__hit_date_iso_bool
\bool_new:N \l__hit_date_digit_bool
\seq_new:N \l__hit_date_seq
\str_new:N \l__hit_date_str
%    \end{macrocode}
%
% 一段是不是纯数字、位数对不对。空串不算。
%    \begin{macrocode}
\cs_new_protected:Npn \__hit_date_digits:nnn #1#2#3
  {
    \bool_set_true:N \l__hit_date_digit_bool
    \bool_lazy_and:nnF
      { \int_compare_p:nNn { \str_count:n {#1} } > { #2 - 1 } }
      { \int_compare_p:nNn { \str_count:n {#1} } < { #3 + 1 } }
      { \bool_set_false:N \l__hit_date_digit_bool }
    \str_map_inline:nn {#1}
      {
        \tl_if_in:nnF { 0123456789 } {##1}
          { \bool_set_false:N \l__hit_date_digit_bool }
      }
  }
%    \end{macrocode}
%
% 拆成年月日。先转成字符串再拆：值里可能有汉字，汉字在 \pkg{xeCJK} 下是活动符，
% 直接拿记号比对会走岔。
%    \begin{macrocode}
\cs_new_protected:Npn \__hit_date_parse:n #1
  {
    \tl_clear:N \l__hit_date_year_tl
    \tl_clear:N \l__hit_date_month_tl
    \tl_clear:N \l__hit_date_day_tl
    \bool_set_true:N \l__hit_date_iso_bool
    \str_set:Nx \l__hit_date_str { \tl_to_str:n {#1} }
    \seq_set_split:NnV \l__hit_date_seq { - } \l__hit_date_str
    \int_compare:nNnTF { \seq_count:N \l__hit_date_seq } > { 3 }
      { \bool_set_false:N \l__hit_date_iso_bool }
      {
        \__hit_date_take:nnnN { 1 } { 4 } { 4 } \l__hit_date_year_tl
        \__hit_date_take:nnnN { 2 } { 1 } { 2 } \l__hit_date_month_tl
        \__hit_date_take:nnnN { 3 } { 1 } { 2 } \l__hit_date_day_tl
      }
  }
\cs_new_protected:Npn \__hit_date_take:nnnN #1#2#3#4
  {
    \int_compare:nNnT { \seq_count:N \l__hit_date_seq } > { #1 - 1 }
      {
        \exp_args:Nx \__hit_date_digits:nnn
          { \seq_item:Nn \l__hit_date_seq {#1} } {#2} {#3}
        \bool_if:NTF \l__hit_date_digit_bool
          { \tl_set:Nx #4 { \int_eval:n { \seq_item:Nn \l__hit_date_seq {#1} } } }
          { \bool_set_false:N \l__hit_date_iso_bool }
      }
  }
%    \end{macrocode}
%
% 月与日经 \cs{int\_eval:n} 走一趟，\texttt{06} 会变成 \texttt{6}：中文排「6 月」，
% 英文取月份名，两边都用不上前导零。年份四位，走一趟不受影响。
%
% 排出来。中文是数字加汉字单位，英文是月份名；\texttt{zh-big} 那一款用汉字数字，
% 目前没有哪一处用它，留着备用。
%    \begin{macrocode}
\cs_new:Npn \__hit_date_zh:n #1
  {
    \l__hit_date_year_tl 年
    \str_if_eq:nnF {#1} { year }
      {
        ~ \l__hit_date_month_tl 月
        \str_if_eq:nnT {#1} { day } { ~ \l__hit_date_day_tl 日 }
      }
  }
\cs_new:Npn \__hit_date_zh_big:n #1
  {
    \zhdigits { \l__hit_date_year_tl } 年
    \str_if_eq:nnF {#1} { year }
      {
        \zhnumber { \l__hit_date_month_tl } 月
        \str_if_eq:nnT {#1} { day }
          { \zhnumber { \l__hit_date_day_tl } 日 }
      }
  }
\cs_new:Npn \__hit_date_en:n #1
  {
    \str_if_eq:nnTF {#1} { year }
      { \l__hit_date_year_tl }
      {
        \hit@month@name:n { \l__hit_date_month_tl }
        \str_if_eq:nnT {#1} { day } { ~ \l__hit_date_day_tl }
        ,~ \l__hit_date_year_tl
      }
  }
\cs_new:Npn \hit@month@name:n #1
  {
    \int_case:nn {#1}
      {
        { 1 } { January }   { 2 }  { February } { 3 }  { March }
        { 4 } { April }     { 5 }  { May }      { 6 }  { June }
        { 7 } { July }      { 8 }  { August }   { 9 }  { September }
        { 10 } { October }  { 11 } { November } { 12 } { December }
      }
  }
%    \end{macrocode}
%
% 有多少排多少：精度不够时按实际有的那一档排，不要排出空的月和日。
%    \begin{macrocode}
\cs_new:Npn \__hit_date_have:
  {
    \tl_if_empty:NTF \l__hit_date_month_tl
      { year }
      { \tl_if_empty:NTF \l__hit_date_day_tl { month } { day } }
  }
%    \end{macrocode}
%
% \begin{macro}{\hit@date@use:nnn}
% \meta{字段名}\marg{样式}\marg{值}。样式取 \texttt{zh}、\texttt{zh-big}
% 或 \texttt{en}；排到哪一级由值自己定。
%
% 空值原样放过：字段没填就是没填，不是格式问题，不该报警告。
%    \begin{macrocode}
\tl_new:N \l__hit_date_raw_tl
\cs_new_protected:Npn \hit@date@use:nnn #1#2#3
  {
    \tl_if_empty:nF {#3}
      {
        \__hit_date_parse:n {#3}
        \bool_if:NTF \l__hit_date_iso_bool
          { \__hit_date_render:nn {#1} {#2} }
          {
            \__hit_date_warn:nnn {#1} { date-not-iso } {#3}
            #3
          }
      }
  }
%    \end{macrocode}
%
% 一个字段一种毛病只报一次。去重的键要把毛病也带上：同一个字段先写得不够细、
% 后来又改成非 ISO，是两回事，只报头一条会把后一条盖掉。
%    \begin{macrocode}
\cs_new_protected:Npn \__hit_date_warn:nnn #1#2#3
  {
    \prop_if_in:NnF \g__hit_warned_options_prop { #1 / #2 }
      {
        \prop_gput:Nnn \g__hit_warned_options_prop { #1 / #2 } { }
        \msg_warning:nnnn { hithesis } {#2} {#1} {#3}
      }
  }
\cs_generate_variant:Nn \hit@date@use:nnn { nnV }
\cs_new_protected:Npn \__hit_date_render:nn #1#2
  {
    \str_case:nn {#2}
      {
        { zh }     { \exp_args:Nx \__hit_date_zh:n     { \__hit_date_have: } }
        { zh-big } { \exp_args:Nx \__hit_date_zh_big:n { \__hit_date_have: } }
        { en }     { \exp_args:Nx \__hit_date_en:n     { \__hit_date_have: } }
      }
  }
%    \end{macrocode}
% \end{macro}
%
% \begin{macro}{\hit@setup@date}
% 两件事要等基础类加载完：把基础类的 \cs{today} 存下来给 \cs{CJKtoday}\oarg{0} 用，
% 再把 \cs{today} 接过来。共享层排在基础类之前，写在顶层会被基础类盖掉。
%    \begin{macrocode}
\cs_new_protected:Npn \hit@setup@date
  {
    \cs_set_eq:NN \CJK@todaysave \today
    \cs_set:Npn \today { \CJK@today }
  }
%    \end{macrocode}
% \end{macro}
%
%    \begin{macrocode}
%</shared-date>
%    \end{macrocode}
%
% \Finale
